\section{Related Works}

Personalized memory retrieval for LLMs has emerged to complement the fixed context window and enable user-specific, context-aware responses. Unlike standard knowledge-based RAG that targets factual data, personal memory retrieval draws on a user’s own history and preferences~\citep{pan2025secom,wu2024longmemeval,xu2025towards} ({e.g.}, retrieving long-term user-AI dialogue context~\citep{jiang2025know,maharana2024evaluating,wu2024longmemeval} and user-user dialogue~\citep{tan2025personabench}) to fill in missing details and tailor responses. Memory-augmented personalized LLM systems combine an LLM with a non-parametric memory to provide relevant background information. We review related works in three areas: (1) query reformulation, (2) index construction, and (3) retrieval frameworks. 

\textbf{Query reformulation.} These methods expand or refine queries to improve memory retrieval. LD-Agent extracts keyphrases for retrieval~\citep{li2025hello}, MemoCue proposes memory-inspired cue query~\citep{zhao2025memocue}, and LQ-TOD generates task-oriented queries~\citep{chen2025llm}. Other approaches leverage LLM-based query expansion, such as LameR~\citep{shen2024retrieval} and MemInsight~\citep{salama2025meminsight}, to enrich the query with contextual attributes.  

\textbf{Index construction.} Prior work has also emphasized structuring user memory into searchable indices. Text-based methods summarize or cluster memories into compact representations, as in MemoryBank~\citep{zhong2024memorybank}, SeCom~\citep{pan2025secom}, and MemGas~\citep{xu2025towards}, while others rely on reflective or hierarchical summaries~\citep{tan2025prospect}. Graph-based approaches instead capture relational structures, exemplified by A-Mem~\citep{xu2025mem}, THEANINE~\citep{ong2024towards}, Mem0~\citep{chhikara2025mem0}, and Zep~\citep{rasmussen2025zep}. Although differing in representation, these methods share a common assumption: retrieval remains static. Most ultimately depend on standard dense retrievers to encode queries and memory keys, applying a uniform retrieval process regardless of uncertainty or task complexity.

\textbf{Retrieval frameworks.} From early keyword-based search~\citep{robertson2009probabilistic} to dense retrievers like DPR~\citep{karpukhin2020dense} and Contriever~\citep{lei2023unsupervised}, retrieval methods aim to rank memory items by semantic similarity. More advanced encoders such as MiniLM~\citep{wang2020minilm}, MPNet~\citep{song2020mpnet}, and BGE~\citep{luo2024bge} improve efficiency and accuracy, and are widely adopted in multi-session dialog retrieval~\citep{xu2021beyond}. However, these frameworks largely adopt a single-process retrieval paradigm, treating all queries as homogeneous regardless of confidence or task complexity.  

\textbf{However, existing methods overlook the dual-process nature of human memory retrieval.}
Most prior works focus on query reformulation, index optimization, or stronger retrievers, yet they implicitly reduce retrieval to a one-shot recognition process (\emph{Familiarity}). This not only neglects the crucial role of deliberate (\emph{Recollection}), but also ignores the need for adaptive switching between the two paths. Our proposed RF-Mem addresses this gap by introducing a recollection retrieval path and a familiarity uncertainty-driven selection that adaptively alternates between fast Familiarity and deeper Recollection retrieval, thereby enabling more personalized memory retrieval.

% \textbf{However, existing methods overlook the adaptive nature of human memory retrieval.} Prior works either optimize query formulation, improve index structures, or design stronger retrievers, but they all assume a fixed retrieval path. This overlooks the dual-process dynamics where recognition (Familiarity) and deliberate reconstruction (Recollection) must be adaptively balanced. Our proposed RF-Mem addresses this gap by introducing an entropy-driven regulation mechanism that dynamically switches between fast Familiarity retrieval and slow Recollection retrieval, thereby achieving scalable, context-sensitive, and personalized memory access.  

% \textbf{Query-reformulate methods} mainly from a text perspective to expand query information for memory retrieval. LD-Agent extract keryphrase for better memory retrieval~\citep{li2025hello}. MemoCue propose a Strategy-Guided Querying for agent proactive memory recall~\citep{zhao2025memocue}. LQ-TOD proposes a query producer to generate queries for retrieval task-oriented dialog~\citep{chen2025llm}. LameR using retrieval augumeted query generation method for better retrieval passage~\citep{shen2024retrieval}. MemInsight uses an LLM to generate query-related attributes to expand the query for memory retrieval~\citep{salama2025meminsight}.

% \textbf{Index-building method} often considers extracting key information from memory data to build an index for better memory retrieval; it can be categorized into two types, the text type and the graph type. In the text type, MemoryBank first proposes a summary method for building a user memory bank with a summary index~\citep {zhong2024memorybank}. SeCom constructs the memory bank at the segment level for topically coherent key building\citep{pan2025secom}. MemGas further improves memory index consolidation by employing Gaussian Mixture Models to construct a multi-granularity association~\citep{xu2025towards}. Reflective Memory Management creates a topic summary-based memory organization by using forward- and backward-looking reflections\cite{tan2025prospect}. In the graph type, A-Mem creates an interconnected knowledge network index in user memory~\citep{xu2025mem}. THEANINE constract user memory graph by considring timelines link~\citep{ong2024towards}. Mem0 pleverages graph-based memory representations to capture complex relational structures among conversational elements~\citep{chhikara2025mem0}, but this method suffers from long user memory~\citep{jiang2025know}. Zep employs episode subgraph, semantic subgraph, and community subgraph to build a user memory index~\citep{rasmussen2025zep}.


% \textbf{Retrieval framework} are proposed for retrieval on user memory, which evolve from early keyword-based search that matches surface tokens to dense embedding methods that capture semantic similarity in continuous vector space. Keyword retrieval provides high precision~\citep{robertson2009probabilistic} on exact matches but struggles with paraphrases or contextual variations. The basic dense retriever~\citep{karpukhin2020dense} used in Multi-Session Chat~\citep{xu2021beyond} to select the top N relevant summaries. Similarly, Contriever~\citep{lei2023unsupervised} are trained using contrastive learning to enhance semantic retrieval. Advanced methods like MiniLM\citep{wang2020minilm}, MPNet~\citep{song2020mpnet}, and BGE~\citep{luo2024bge} have been proposed using a pre-training method for efficient retrieval. 